\subsection{A complete support criterion for nonlinear Hopf stars}

\begin{statement}{C.2.UN.49}{Proved}{Exact nonlinear Hopf star criterion and compact realization}
Let $S$ be a specified compact SNC surface with components
$P,H_1,\ldots,H_n$, $n\ge1$, and no triple points. Suppose
$q:P\to B$ is a connected fixed-Jacobian elliptic quasi-bundle:
its only nonsmooth scheme fibres are multiples of smooth
elliptic curves, and its integral $H^1$ monodromy is the
identity. Each $H_i=H_{r_i}(\beta_i)$ is a nonlinear primary
Hopf surface as in (C.NH.1). It meets $P$ along its unique
axis, identified by a specified holomorphic isomorphism
with a distinct reduced fibre $J_i$ of $q$. There are no
other intersections.

The following assertions are equivalent.

\begin{enumerate}
\item $S$ is the reduced full central fibre of a proper
one-parameter degeneration of two-tori with smooth total space.
\item $S$ is the reduced full special fibre of a smooth compact
threefold over $\PP^1$, all of whose other fibres are smooth
two-tori.
\item The normal lines on every double curve are dual;
$B$ is elliptic; all $r_i$ equal one integer $r$; the full
marked fibres are $rJ_i$ and there are no other multiple
fibres; and the actual inverse logarithmic transform
$T_0\to B$ defined in (C.NS.1)--(C.NS.3) has zero integral
Chern vector in $H^2(B,\Lambda_F)$.
\end{enumerate}

Here $F$ is the fixed elliptic Jacobian and $\Lambda_F$ its
integral period lattice. The inverse transform retains
the entire analytic torsor of $P$ and the given axis
identifications. It is unique up to the explicit product
translations described below. Thus the criterion concerns
the specified surface $S$, including its component
moduli and gluing maps. Every fibre support in assertion
(2) has Euler number zero. No extension of the original
smoothing's elliptic quotient is assumed.
\end{statement}

\begin{proof}
\par\medskip\noindent\textbf{The actual inverse transform.}
We first define the operation whenever the marked
multiplicity is $r_i$ and the two normal lines are dual.
Write the Hopf axis as $\C^*/\beta_i^{r_i\Z}$, with
coordinate $z$. Its normal has multiplier $\beta_i$.
The matching neighborhood in $P$ has the precise model
\[
 p=x^{r_i}z,\qquad
 (z,x)\sim(\beta_i^{r_i}z,\beta_i^{-1}x).
                                                        \tag{C.NS.1}
\]
Indeed, normalize the $r_i$-th root of $q$ at this fibre.
The result is a smooth principal $F$-bundle over a disc,
hence is trivial there. The free root action is an
order-$r_i$ translation on $F$, together with the
primitive normal character. A finite averaging gauge
makes that translation independent of the root
parameter. The dual-normal isomorphism identifies this
connected cyclic cover with the one coming from the
Hopf axis. In multiplicative coordinates the cover is
$z=w^{r_i}$, $x=u/w$, with deck $w\mapsto\beta_i w$;
its quotient is exactly (C.NS.1), with $p=u^{r_i}$.
This also proves that the normal of a smooth
multiple fibre has exact order its multiplicity.

On the punctured base disc the map
\[
       (z,x)\longmapsto(p,w)=(x^{r_i}z,x^{-1}),
       \qquad z=pw^{r_i},
       \qquad [w]\in F_i=\C^*/\beta_i^{\Z}
                                                        \tag{C.NS.2}
\]
is an actual product trivialization. Replace (C.NS.1)
by $F_i\times\Delta_p$ through (C.NS.2), at every
marked point. The original torsor on the complement
and these exact transition maps define
\[
                             q_0:T_0\longrightarrow B.
                                                        \tag{C.NS.3}
\]
When there are no unmarked multiple fibres, this is a
smooth principal $F$-bundle. All identifications of
$F_i$ with $F$ are the ones furnished by the given
normal covers and the original torsor.

The construction is independent of the permitted local
choices. Two dual-normal isomorphisms differ by a
nonzero holomorphic function on the compact elliptic
axis and hence by a constant. This changes $w$ by a
constant translation of $F_i$, which extends over the
product disc. An automorphism of a root-trivialized
principal $F$-bundle that extends across its central
curve is a holomorphic translation there, followed by
a fixed linear automorphism. If it preserves the
specified normal cover, that translation is invariant
under the root action and descends to a holomorphic
translation on the ordinary disc. It again extends
over the product filling. The specified axis map
fixes the remaining linear identification. Thus
(C.NS.3) retains the full holomorphic torsor class;
it is not defined merely by a rational Chern number.

\par\medskip\noindent\textbf{The integral boundary matrix and the smoothing topology.}
We compute all meridians before applying a smoothing
argument. In (C.NS.1), let $a$ be the positive angular
loop of $z$, let $b$ be its deck loop, and let
$\mu_P$ be the positive angular loop of $x$. On the
punctured product (C.NS.2), let $\delta$ be the
positive base loop with $w$ fixed, $e$ the angular
loop of $w$, and $h$ its $\beta_i$-deck loop. Then
\[
 \begin{pmatrix}a&b&\mu_P\end{pmatrix}_{(\delta,e,h)}
       =
 \begin{pmatrix}
  1&0&r_i\\
  0&0&-1\\
  0&1&0
 \end{pmatrix};
 \qquad
 a=\delta,\quad b=h,\quad \mu_P=r_i\delta-e.
                                                        \tag{C.NS.4}
\]
The matrix is integral of determinant one. For example,
holding $p$ fixed while $w$ winds once gives
$x=w^{-1}$ and $z=pw^{r_i}$, proving
$e=r_i a-\mu_P$. The formulas retain the normal
character, not just its order.

The nonlinear Hopf pair is smoothly equivalent, for
this boundary calculation, to the diagonal pair with
contraction $(\beta_i^{r_i}z,\beta_i w)$. One exact
comparison on the complement is
\[
 U={z\over w^{r_i}},\qquad
 V=U-{\beta_i^{-r_i}\over\log|\beta_i|}\log|w|.
                                                        \tag{C.NS.5}
\]
The deck map sends $U$ to $U+\beta_i^{-r_i}$ and
$w$ to $\beta_iw$, so $V$ is invariant. This gives
$H_i\setminus J_i\simeq\C_V\times F_i$ smoothly.
The map
$z\mapsto z-\beta_i^{-r_i}w^{r_i}\log|w|/\log|\beta_i|$
extends continuously as the identity on the axis and
preserves orientation; it is smooth on every boundary
torus. Equivalently, vary the resonant coefficient from
one to zero in the contraction and use a uniform
compact fundamental annulus. This is a smooth proper
family of pairs, so the same integral boundary
identification follows by relative Ehresmann.

Let $\mu_H$ be the Hopf normal meridian. The
smoothing equation is $xw=t$, so
$\mu_H=-\mu_P$. The regular fibre cycles on the
Hopf side are
\[
                   e_H=r_i a+\mu_H=e,\qquad h_H=b=h.
                                                        \tag{C.NS.6}
\]
The complementary solid torus $F_i\times D^2$ in
(C.NS.5) kills $a=\delta$. Consequently replacing
the tubular neighborhood in $P$ by the Hopf
complement is precisely the inverse logarithmic
transform (C.NS.2), with its indicated integral
framing.

For completeness, this calculation applies to every
actual SNC smoothing of the specified support.
Remove small tubular neighborhoods of the disjoint
double curves. On each remaining compact component
core the original map is a submersion through zero;
Ehresmann with boundary identifies that core with
its nearby copy. On the collars the local equations
are exactly $xw=t$ after absorbing a holomorphic
unit. The product of the two normal lines is
trivialized by this same function. Two such
trivializations along a compact double curve differ
by a constant in $\C^*$, whose rotation is isotopic
to the identity. Thus the nearby four-manifold is
the indicated tubular plumbing, and (C.NS.4)--(C.NS.6)
identify it with the underlying four-manifold of
$T_0$. Under this identification, the oriented
generator of $H_2(H_i\setminus J_i,\Z)$ is the
ordinary fibre class of $q_0$, for every $i$.
The identifications between the normal covers are
holomorphic elliptic isomorphisms, hence have degree
one and preserve that orientation. In particular all
these transported core cycles are the same integral
two-cycle in the nearby fibre, even when their
primitive angular directions in $\Lambda_F$ differ.

Here is an explicit global version of the Chern test.
Trivialize the principal $F$-bundle on the punctured
base and compare that trivialization with (C.NS.2)
around the $i$-th disc. Let $\nu_i\in\Lambda_F$ be
the winding of its translation transition, with
positive disc-boundary orientation. Then
\[
 c(T_0)=\sum_i\nu_i,\qquad
 \mu_{P,i}=r_i\delta_{\rm ext}
                    +(r_i\nu_i-e_i).
                                                        \tag{C.NS.7}
\]
The first equality is the integral obstruction to
extending these boundary trivializations over the
discs. It can alternatively be used as the sign
convention for the Chern vector. The second equality
is (C.NS.4) after $\delta_i=\delta_{\rm ext}+\nu_i$.
Thus the rational Seifert vector in this convention
is $\sum_i(\nu_i-e_i/r_i)$, and the inverse transform
adds exactly $\sum_i e_i/r_i$. The holomorphic
transition functions themselves, rather than only
their windings, remain in (C.NS.3).

\par\medskip\noindent\textbf{Canonical orders, the elliptic base, and equal resonance.}
Suppose first that $S$ occurs as in assertion (1).
The equation $xw=t$ proves duality of the normal
lines. Their exact torsion orders then give the
marked multiplicities $r_i$.

Choose a generator of the rank-one direct image of
the relative canonical line over the disc. Its
evaluation has a vertical divisor
$k_PP+\sum_i k_iH_i$. It is nowhere zero on every
nonzero torus fibre. Adjunction on the Hopf component
and $P|_{H_i}=J_i$, $H_i|_{H_i}=-J_i$ give
\[
 (k_P-k_i)J_i=K_{H_i}+J_i=-r_iJ_i,
       \qquad k_i-k_P=r_i.
                                                        \tag{C.NS.8}
\]
The last equality is an equality of integers:
$\OO_{H_i}(J_i)$ has infinite order, since a
positive effective divisor on a compact complex
surface cannot have a trivial positive tensor power.
Adjunction on $P$ now gives the actual line equality
\[
                         K_P=\OO_P\Bigl(\sum_i(r_i-1)J_i\Bigr).
                                                        \tag{C.NS.9}
\]

For the given quasi-bundle the root charts give
$q_*K_P=K_B$ and
$K_P=q^*K_B+\sum_{\text{all multiple fibres}}(m-1)J$.
Indeed, the local form is the base differential
wedged with the invariant elliptic differential;
on the root chart its exact zero is $(m-1)J$,
and taking the fibrewise constant Laurent
coefficients gives the displayed pushforward.
Equation (C.NS.9) implies $h^0(K_P)\ge1$, so
$g(B)\ge1$. Subtracting the marked terms from
the two canonical formulas gives
\[
 q^*K_B+\sum_{\text{unmarked multiple fibres}}(m-1)J=0
                    \quad\hbox{as a line bundle}.       \tag{C.NS.10}
\]
For $g(B)\ge2$ a nonzero canonical section on $B$
has a nonzero effective divisor, and its pullback
together with the other terms would be a nonzero
effective divisor linearly equivalent to zero on
the compact surface. This is impossible. Therefore
$B$ is elliptic, and the same argument in (C.NS.10)
removes every unmarked multiple fibre.

Let $\omega_t$ be the evaluated relative two-form
divided by $t^{k_P}$. On a compact core of $H_i\setminus J_i$
its leading term is $t^{r_i}\eta_i$, where
\[
 \eta_i=c_i\,{dz\wedge dw\over w^{r_i+1}},
                         \qquad c_i\ne0.                \tag{C.NS.11}
\]
To verify this exact pole order, write the total
canonical evaluation in a double chart as
$x^{k_i}w^{k_P}$ times a holomorphic unit and a
coordinate three-form, then divide by $dt$ with
$t=xw$. Equation (C.NS.8) gives (C.NS.11).
The resulting coefficient on the entire Hopf
component is a nowhere-zero section of
$K_{H_i}((r_i+1)J_i)=\OO_{H_i}$, so it is the
nonzero constant $c_i$.

On the logarithmic cover of $H_i\setminus J_i$,
put $U=z/w^{r_i}$ and $W=\log w$. Its two deck
periods are
$(0,2\pi i)$ and $(\beta_i^{-r_i},\log\beta_i)$.
Thus the integral of $\eta_i=c_i\,dU\wedge dW$
on its oriented core torus is
$\pm2\pi i c_i\beta_i^{-r_i}\ne0$.
Transport that compact torus through the
submersion near the core. By (C.NS.6) all such
cycles are the same class $\gamma$ in the
nearby fibre. Therefore the single holomorphic
period satisfies, for every $i$,
\[
        \int_\gamma\omega_t
             =t^{r_i}\bigl(\pm2\pi i c_i\beta_i^{-r_i}
                              +O(t)\bigr).              \tag{C.NS.12}
\]
Its vanishing order is simultaneously every
$r_i$. Hence all resonances equal $r$.
This uses the transported two-cycle; it makes
no assumption about parallel primitive angular
directions or an elliptic quotient of the
nearby complex torus.

By the plumbing comparison the underlying
four-manifold of $T_0$ is a real four-torus.
For a principal elliptic bundle over the
elliptic curve, the homotopy sequence injects
$\Lambda_F$ into its fundamental group, and
the exact presentation is
\[
 \langle a,b,e,h:\ e,h\text{ central},\
             [a,b]=e^{c_1}h^{c_2}\rangle,
                         \quad(c_1,c_2)=c(T_0).          \tag{C.NS.13}
\]
Its fundamental group is abelian only if
$c(T_0)=0$. This proves every necessity in (3).

\par\medskip\noindent\textbf{The inverse transform is then an actual torus.}
Assume (3). The sheaf sequence
$0\to\Lambda_F\to\OO_B\to F\to0$ lifts the
principal torsor class of $T_0$ to $H^1(B,\OO_B)$
because its connecting class is precisely
$c(T_0)=0$. On $B=\C/\Gamma_B$ represent this
class by a constant $(0,1)$-form. Integrating it
on the universal cover gives the actual
translation lattice
\[
 (p,z)\mapsto(p+\gamma,z+a\bar\gamma),\quad
 (p,z)\mapsto(p,z+\lambda),
       \qquad \gamma\in\Gamma_B,\quad\lambda\in\Lambda_F.
                                                        \tag{C.NS.14}
\]
Changes of the chosen Cech cocycle are fibre
translations and give an isomorphism with the
original torsor, exactly as in (C.KX.14).
The lattice is discrete of real rank four.
Consequently $T_0$ is a complex two-torus,
with its specified principal elliptic quotient
and all the marked fibres retained.

\par\medskip\noindent\textbf{A compact twistor family with the exact local period presentations.}
Write $T_0=\C^2/\Lambda$ in coordinates $(p,z)$
so that the common elliptic fibre direction is
$p=0$. On the affine $t$-chart of $\PP^1$ put
\[
 \ell_t(p,z)=(p+t\bar z,\ z-t\bar p).
                                                        \tag{C.NS.15}
\]
This is real-linear and invertible for every
$t$, since its inverse has denominator
$1+|t|^2$. The values $\ell_t(v)$ for
$v\in\Lambda$ are holomorphic sections of
$\OO_{\PP^1}(1)^{\oplus2}$: on the other chart
they are $(s p+\bar z,s z-\bar p)$, $s=t^{-1}$.
Translate the vector bundle by these sections
and take the quotient. Each fibre lattice has
real rank four. The inverse real-linear maps
give local uniform fundamental parallelepipeds
over the compact base, proving a free proper
action and a smooth compact quotient
$\mathcal T\to\PP^1$. Its fibres are two-tori
and its fibre at zero is the actual $T_0$.
Pull this family back by $t=s^r$; the pullback
remains smooth and compact because the
original map is a submersion.

For the $i$-th Hopf identification choose the
oriented fibre-lattice basis
$\lambda_i=(0,a_i)$, $\mu_i=(0,b_i)$ so that
$\beta_i=\exp(2\pi i b_i/a_i)$. This is the
basis supplied by (C.NS.2), not a replacement
of the given Hopf modulus. On the vector
space with coordinates $(P,Z)$ after
(C.NS.15), define
\[
 p_i=P-s^r{\bar a_i\over a_i}Z,\qquad
 z_i={2\pi i\over a_i}Z.
                                                        \tag{C.NS.16}
\]
The $\lambda_i$-period is $(0,2\pi i)$ and
the $\mu_i$-period is
\[
   \bigl(\delta_i(s),\log\beta_i\bigr),\qquad
   \delta_i(s)=s^r
        \left(\bar b_i-{\bar a_i b_i\over a_i}\right).
                                                        \tag{C.NS.17}
\]
The coefficient is nonzero because $a_i,b_i$
are an oriented real basis of a complex line.

Complete this pair to an integral basis of
$\Lambda$ by two lifts of the base lattice.
Their $p_i$-periods define a smooth local
family of elliptic curves $B_{i,s}$, and
their exponentiated $z_i$-periods define a
holomorphic family of flat degree-zero line
bundles $M_i$ on $B_{i,s}$. Thus, on a
sufficiently small disc, the actual torus
family has the exact presentation
\[
 \mathcal T_s
   =M_{i,s}^{\times}/\langle g_i\rangle,\qquad
 g_i(p,\ell)=(p+\delta_i(s),\beta_i\ell).
                                                        \tag{C.NS.18}
\]
This equality follows by quotienting the
same four displayed lattice generators,
first by the three excluding $\mu_i$.
It retains the flat part of the torsor,
even when that part is not torsion.

\par\medskip\noindent\textbf{The local weighted insertion with the varying flat line.}
At the marked point of $B_{i,0}$ perform
the weight-$(r,1)$ blowup of the relative
elliptic surface $B_i\to\Delta_s$, and replace
$M_i$ by its pullback tensored with
$\OO(R_i)$. The nonzero line has the same
smooth coarse chart as (C.NH.6)--(C.NH.7).
If $\delta_i(s)=d_i s^r$, the lifted deck
map in a flat local frame is exactly
\[
 (z,w,x)\mapsto
   \bigl(\beta_i^r(z+d_iw^r),
              \beta_iw,\beta_i^{-1}x\bigr),
           \qquad p-p_i=x^rz,\quad s=xw.
                                                        \tag{C.NS.19}
\]
The change $Z=z/(\beta_i^r d_i)$ gives
the specified $H_r(\beta_i)$ on $x=0$.
On $w=0$ the chart is precisely (C.NS.1);
hence the other component is obtained
from $T_0$ by the inverse of (C.NS.2).

We check properness with the varying
flat line. On a small closed base disc,
choose smoothly varying unitary flat
metrics on $M_{i,s}$. If its period
characters are $e^{Z_j(s)}$ and base
periods $P_j(s)$, solve the two real
equations
$\operatorname{Re}(Z_j-c(s)P_j)=0$.
Their determinant is nonzero, so $c(s)$
is smooth and bounded. In the resulting
metric, the norm multiplier of the flat
lift in (C.NS.18) is
$|\beta_i|\exp(-\operatorname{Re}(c(s)\delta_i(s)))$.
It is uniformly less than one after
shrinking the disc. The additional
exceptional line has the two-sided
polynomial distortion (C.NH.9)--(C.NH.10).
Exponential contraction still dominates
that polynomial uniformly on the closed
disc. The compact-annulus proof of
(C.NH.11) therefore gives a smooth
proper filling $\mathcal Y_i\to\Delta_s$.
Its infinite cyclic action is free,
and every prescribed $\beta_i$ is
retained.

Off the exceptional divisor the canonical
section of $\OO(R_i)$ identifies the new
line with $M_i$, equivariantly for the
two deck maps. Consequently there is
the actual isomorphism
\[
       \mathcal Y_i\setminus H_i
            \simeq \mathcal T|_{\Delta_s}\setminus F_i,
                                                        \tag{C.NS.20}
\]
where $F_i$ is the marked central
elliptic fibre in $T_0$. On the remaining
central component this is exactly the
punctured product map (C.NS.2).

\par\medskip\noindent\textbf{Simultaneous replacement and compactness.}
The central curves $F_i$ are disjoint.
We justify that (C.NS.20) allows
independent replacements near them,
although its generic tori need not
have an elliptic quotient.

If a sequence approaches $H_i$ in
$\mathcal Y_i$, use a compact Hopf
fundamental domain in (C.NS.19).
Then $z,w$ stay bounded, do not
simultaneously approach zero, and
$x\to0$. Its image in the original
line presentation has
$p-p_i=x^rz$, $s=xw$, and flat
line coordinate proportional to $x^{-1}$.
To put that coordinate in a fixed
compact norm annulus requires
$O(1+|\log|x||)$ deck iterates.
Their total base displacement is
$O(|s|^r(1+|\log|x||))$, which
tends to zero. Hence every cluster
point of the images lies in $F_i$.
Conversely properness of $\mathcal Y_i$
shows that images approaching a point
outside $F_i$ have only their unique
inverse-image limit outside $H_i$.

Choose disjoint relatively compact
neighborhoods $W_i$ of $F_i$ in
$\mathcal T$ over a small base disc.
The set consisting of $H_i$ and the
inverse image of $W_i\setminus F_i$
under (C.NS.20) is open, by the
cluster-set calculation. Glue these
open replacements to
$\mathcal T\setminus\bigcup_i F_i$.
Their identification graphs are
relatively closed by that same
calculation, and the replacements
have disjoint central neighborhoods.
Thus the finite gluing is Hausdorff
and smooth. Nested compact
neighborhoods inside the $W_i$,
properness of each $\mathcal Y_i$,
and the compact original complement
cover every closed inverse image
of a sufficiently small base disc.
The resulting map is proper.

The central main component is exactly
$P$: outside the marked fibres its
torsor is unchanged, and at every
marked fibre (C.NS.19) reverses
the precise map (C.NS.2).
The Hopf components are the prescribed
$H_i$. Any constant axis translation
in their identification is realized
by the Hopf automorphism
$(z,w)\mapsto(a^rz,aw)$; the linear
part has already been retained in
the chosen fibre-lattice basis.
Thus all specified elliptic gluing
maps, not just their degrees, are
the original ones. The full central
equations are $xw=s$, so the fibre
is reduced SNC.

Finally carry out these replacements
inside the compact pulled-back
twistor family over $\PP^1$ and
leave its complement unchanged.
The preceding proper local gluing
and the original compact complement
give a smooth compact total space.
Every fibre away from zero remains
the original two-torus. The central
support has Euler number zero since
each elliptic quasi-bundle, Hopf
component, and double elliptic curve
has Euler number zero. We have
proved (3)$\Rightarrow$(2), including
the full compact-total-space and
torsor comparisons. Assertion
(2)$\Rightarrow$(1) is restriction
to a disc, completing the equivalence.
\end{proof}
